\chapter{Robust, Flexible and Imputation Methods}

\section{Robust M-estimation}

`rlm` fits a linear model with M-estimators that down-weight outliers.
The default is a Huber loss with a tuning constant, which gives OLS-like
behaviour near the centre and L1-like behaviour in the tails.

\begin{lstlisting}[caption={Robust regression}]
let m_rlm = rlm(y ~ x1 + x2, df)
print(m_rlm)
\end{lstlisting}

Use `rlm` when the data contain influential outliers but the majority of
the distribution is approximately normal.

\section{Generalised Linear Models}

`glm` fits a linear predictor through a link function. It generalises
OLS to exponential families such as Poisson, Binomial, and Gamma.

\begin{lstlisting}[caption={GLM}]
let m_glm = glm(y ~ x1 + x2, df, family=poisson, link=log)
print(m_glm)
\end{lstlisting}

Available families include `gaussian`, `poisson`, `binomial`, `gamma`,
`inverse_gaussian`, and `negative_binomial`. Links include `identity`,
`log`, `logit`, `probit`, `cloglog`, and `inverse`.

\section{GEE}

Generalised Estimating Equations fit a marginal model while accounting
for within-group correlation. It is the population-averaged analogue of
mixed models.

\begin{lstlisting}[caption={GEE}]
let m_gee = gee(y ~ x1 + x2, df, id="group", family=poisson)
print(m_gee)
\end{lstlisting}

Correlation structures include `independence`, `exchangeable`, `ar1`, and
`unstructured`.

\section{Beta regression}

`betareg` is appropriate when the dependent variable is bounded in
$(0,1)$.

\begin{lstlisting}[caption={Beta regression}]
let m_beta = betareg(share ~ x1 + x2, df)
print(m_beta)
\end{lstlisting}

It models the mean through a logit link and the dispersion through a
separate precision parameter.

\section{GAM and lowess}

`gam` fits a generalised additive model with smoothing splines for
nonlinear effects.

\begin{lstlisting}[caption={GAM}]
let m_gam = gam(y ~ x1 + s(x2) + x3, df)
print(m_gam)
\end{lstlisting}

`lowess` is a local polynomial smoother for visual exploration and
partial regression.

\begin{lstlisting}[caption={Lowess}]
let m_low = lowess(df, y, x2, frac=0.3)
print(m_low)
\end{lstlisting}

\section{Multiple imputation}

`mice` performs multiple imputation by chained equations. It fills in
missing values using a sequence of conditional models and returns a set
of completed datasets.

\begin{lstlisting}[caption={MICE}]
let m_mice = mice(df, vars=["y", "x1", "x2"], m=5)
print(m_mice)
\end{lstlisting}

Rubin's rules combine the results from each imputed dataset into a
single estimate with a proper variance that reflects imputation
uncertainty.
